% Generated draft: requires mathematical and source review.
\documentclass[11pt]{article}
\usepackage[margin=25mm]{geometry}
\usepackage{fontspec}
\setmainfont{Latin Modern Roman}
\usepackage{amsmath,amssymb,mathtools,unicode-math}
\setmathfont{Latin Modern Math}
\usepackage{enumitem,needspace,xurl,hyperref}
\setlength{\parindent}{0pt}
\setlength{\parskip}{5pt}
\setlength{\emergencystretch}{3em}
\newcommand{\N}{\mathbb{N}}
\newcommand{\Z}{\mathbb{Z}}
\newcommand{\Q}{\mathbb{Q}}
\newcommand{\R}{\mathbb{R}}
\newcommand{\C}{\mathbb{C}}
\newcommand{\F}{\mathbb{F}}
\newcommand{\A}{\mathbb{A}}
\newcommand{\PP}{\mathbb{P}}
\newcommand{\E}{\mathbb{E}}
\newcommand{\eps}{\varepsilon}
\newcommand{\dd}{\,\mathrm{d}}
\begin{document}
\section*{UnsolvedMath ID: 20003287}
\textbf{Source problem code:} catalogue entry
\subsection{Definitions and mathematical statement}
The source contains two distinct items. Conjecture 18, attributed to N. Shah, says only that ``the same result'' holds for every Lie group $L$ containing $G=SO(n,1)$. The antecedent result is absent from this record, so the conjecture is not a standalone precise theorem here. The aid reports its resolution by Lei Yang's Theorem 1.4 under the exact lattice, dense-orbit, analyticity and visual-nondegeneracy hypotheses. It does not reproduce those hypotheses, the curve family, target measure or conclusion. Accordingly the reported solved status is retained without inventing the missing expanding-curve statement or asserting it for all curves or all orbits.

The next paragraph is identified by the aid as setup for Question 19, not a clause of Conjecture 18. Independently let a semisimple Lie group $G$ act linearly on a finite-dimensional normed real vector space $V$, fix $p\in V$, and put
\[
 R_T=\{g\in G:\|gp\|<T\},\qquad T>0.
\]
For a lattice $\Gamma<G$ and quotient map $\pi:G\to G/\Gamma$, the proposed measure is
\[
 \mu_T=\pi_*\left(\frac{\mathbf1_{R_T}\,dg}{\operatorname{vol}_G(R_T)}\right),
\]
which requires $0<\operatorname{vol}_G(R_T)<\infty$. The source's assertion that finite stabilizer makes every $R_T$ compact is false as stated. Noncompactness alone, however, does not automatically imply infinite Haar volume.

The aid's explicit correction uses $G=SL(2,\R)$ acting on binary cubics by the row-variable convention $(g\cdot f)(x,y)=f((x,y)g)$, with $p=x^2y$. The stabilizer of $p$ is reported trivial, but
\[
 a_s=\operatorname{diag}(e^{-s},e^s),\qquad
 a_s\cdot p=e^{-s}p\longrightarrow0
 \quad(s\to\infty).
\]
The sequence $a_s$ escapes compact subsets while eventually lying in $R_T$ for every $T>0$. Thus these strict sublevels are not even relatively compact, and the orbit is not closed because $0$ is not in it. The obstruction is properness of the orbit map, not merely stabilizer size.

In the standard finite-dimensional Lie-group representation framework, the aid reports the precise criterion
\[
 g\mapsto gp\text{ is proper}
 \quad\Longleftrightarrow\quad
 G_p\text{ is compact and }G\cdot p\text{ is closed}.
\]
With this property, the \emph{closed} norm sublevels $\{g:\|gp\|\leq T\}$ are compact and the strict $R_T$ are relatively compact. Strict sublevels are generally open, so they are not silently called compact. Finite stabilizer supplies compactness of $G_p$ only; it does not supply the closed-orbit condition. Proper sublevels ensure finite Haar mass, with nonemptiness still needed for normalization. This repairs the Question 19 setup but does not state or solve its omitted equidistribution question, nor replace the missing hypotheses of Yang's theorem for Conjecture 18.
\subsection{Short English statement}
The reported resolution of Shah’s conjecture cannot be fully restated without its missing antecedent and theorem hypotheses. The adjacent representation paragraph has a separate error: a trivial stabilizer need not make norm sublevels relatively compact. The aid supplies an explicit binary-cubic example and the proper-orbit correction.
\subsection{Sources}
\begin{itemize}
\item Local UnsolvedMath registry, permanent ID 20003287, source code catalogue entry; \texttt{aim-other-notes.json}, notes[93], source tag conjecture.
\item AIM, \emph{Open Problems from the Workshop Emerging Applications of Measure Rigidity}, Section 4.3, Conjecture 18, N. Shah, followed by the setup for Question 19: \url{https://aimath.org/WWN/measrigid/measrigid.pdf}.
\item Lei Yang, Theorem 1.4, cited by the research summary for the conjecture's resolution; no title, date, URL or full theorem statement is supplied.
\item The binary-cubic counterexample and proper-orbit criterion are supplied by the aid without separate full references.
\item Machine-generated research attempt 1: partial result, solved status at the run, rigor labelled ``proved,'' low novelty confidence; registry solved/``SOLVED-IN-LITERATURE.'' Not peer reviewed; independently verify before citation. Registry updated 14 August 2026.
\end{itemize}
\end{document}
